% Page budget: 1.55 pages | Owns: augmentation topology, additional states, encoder, and closed-loop objective; per README ownership also the training coordinates of theta_base, the optimiser, checkpoint selection and the added-state initialisation (PS2); models M_U and M_PS2.
% WRITING GUIDE
% Purpose: Define the final augmentation architecture and explain how it is trained inside the known feedback loop.
% Start here: Current implementation decisions D-129, D-140, D-141, D-180, D-220, D-233, and D-237 in docs/decisions.md.
% Supporting material: Quinten's 18 September outline feedback, the Meeting-audit synthesis, Thesis-writeup/Documentation/TRAINING-DESIGN.md, and docs/writeup figures.
% Mandatory papers: Read Hoekstra's 2025 LFR augmentation paper, the 2026 first-principles augmentation paper, the encoder-initialization paper, and Kessels' Chapter 5 before writing this section.
% Research-plan anchor: Preserve Aspect 2's dynamic parallel augmentation objective; document the departure from open-loop training and earlier state routing.
% Current decisions: Separate physical and additional states, keep the controller outside the model, score the complete training window without burn-in, and use one network for the physical correction and added-state update.
% Choices to justify: Final reported routing, added-state dimension, ANN width and depth, zero initialization, encoder history, closed-loop objective, rollout length, full-window scoring, and validation horizon. Use the configuration of the reported thesis run, not a superseded routing experiment.
% Horizon check: Treat encoder history, scored training window, joint-estimation horizon, and free-run validation horizon separately; relate candidates to relevant stable modes and test sensitivity to slower dynamics and free-integrator accumulation.
% Implementation audit: Read scripts/gantry/gantry_interconnect_dynamic.py, gantry_dynamic/{model,config,controller,data,training}.py, and model_augmentation/fit_systems/{interconnect,closed_loop,pre_encoder}.py.
% Reference check: Verify sources for parallel LFR augmentation, encoder initialization, closed-loop state extension, and any claimed architectural precedent against the original paper or thesis.
% Cautions: Earlier routing plans are superseded; distinguish the truth-only hidden absorber from learned states; closed-loop fit alone does not prove autonomous plant recovery.
% Figures: Absorber schematic, author's updated architecture, and thesis-owned common-reference loops. No residual panel or horizon figure.
% Section is finished when: Every signal has a defined frame and timing, the RK4 boundary is explicit, and the described architecture matches the final code.
%
% SOURCE MAP
% Hoekstra et al. 2025 EJC: dynamic-parallel topology, additional states, encoder-based truncated simulation.
% Hoekstra et al. 2026 Automatica submission: extended LFR formulation, baseline-preserving model/encoder initialization, joint identification context.
% Hoekstra et al. 2026 encoder paper: reconstructability-based W^b initialization; it does not identify W^a or test dynamic augmentation.
% Kessels 2025 thesis: direct closed-loop rollout, controller equations, and reconstruction of the machine controller state for each window.
% Drenth 2025 thesis: LPV-specific augmentation precedent; not the unrestricted MLP structure used here.
% Hefny et al. 2015, Downey et al. 2017: two-stage regression and its use as an RNN initialisation before BPTT (PS2 motivation only).
% This project: gantry specialization, residual-loop implementation, exact routing, RK4 boundary, PS2, and verification.
%
% MUST ESTABLISH (agreed with Dirk, 2026-10-06; step 2a of the thesis-section skill)
% Contribution delivered: the gantry realisation of S-DP augmentation (CT self-scheduled LPV-LFR baseline
%   inside a DT parallel model, joint estimation in the ten combinations) and its identification under
%   feedback. Not new: the S-DP topology (Hoekstra), augmentation of self-scheduled LPV-LFR with added
%   states (Drenth Eq. 5.2), closed-loop training with a known controller (Kessels, inspiration only).
% III-A Augmented model
%  1. Dynamic, not static: static augmentation adds no states (EJC Sec. 2, "flexible modes"); refitting
%     theta_base adds none either. Argue from the general case, no absorber.
%  1b. The augmented model is Hoekstra's LFR augmentation structure (EJC Eq. 4, 2026 Eq. 6) with a FIXED
%     interconnection (EJC Sec. 3.2) realising S-DP; its baseline block is Section II's self-scheduled
%     LFR with Delta(Y) = Y I6: an LFR inside an LFR, which is what the section title means.
%     Well-posedness: acyclic graph (EJC Thm. 1, 2026 Cond. 6) plus Section II's M(Y) condition.
%  2. Parallel, not series. LEAD REASON (Dirk, 2026-10-06): the orthogonal-by-construction method that
%     Section IV builds on is defined for the additive (parallel) structure (Gyorok 2026 abstract, Sec. 1;
%     D-003), so parallel lets us use that existing work. Supporting: Hoekstra gives no universal rule,
%     the structure follows prior knowledge (2026 Sec. 1, 6.3); the baseline stays the explicit nominal
%     term (research plan Aspect 2); parallel needs only a feedforward network for baseline-preserving
%     init, series needs ResNets (EJC Sec. 4.3, 2026 Sec. 5.4, 6.3).
%  3. State augmentation only, h_aug = 0: output P^T q is exact kinematics (cf. Kessels Remark 5.3). TODO reason.
%  4. Correction on all physical rows and added states: S-DP definition (EJC), X and Y reachable (D-103);
%     marginal-mode risk on free axes stated. TODO: why correct position rows rather than keep kinematics.
%     Context: Hoekstra 2026 Sec. 7 (F1Tenth) detaches the free integrators and learns only the
%     input-to-velocity map because the full state is measured; here only positions are measured.
%     One network for f_aug and g_aug is an implementation statement (D-180), no scientific claim.
%  5. DT correction after the CT RK4 step (input held, M(Y) per stage, affine normalisation). TODO reason
%     versus force injection in the CT derivative.
%  6. Unrestricted nonlinear augmentation (depends on Y through the state) versus LPV-structured
%     augmentation (Drenth): departure from the research plan's "in LPV-LFR form". TODO reason. Acyclic
%     routing adds no instantaneous loops; baseline well-posedness is Section II's.
%  7. Zero output layer reproduces the baseline from the same physical initial state (Hoekstra 2026
%     Eq. 29, D-237); the split is not unique (Sec. 5.2), handed to Section IV. Added states are memory
%     and model order, defined up to a state transformation, no physical coordinates.
% III-B Encoder
%  8. Encoder needed: only positions measured; truncated windows for cost and stability (Beintema).
%  9. W^b from the reconstructability map of a ZOH linearisation at an equilibrium (encoder paper
%     Eqs. 15-17, 30-32), not fitted (Hoekstra 2026 Eq. 30); W^a random. Claim only better starting
%     values and faster convergence, final accuracy comparable (encoder paper Sec. 4.4). The paper
%     assumes a stable baseline; ours is marginally stable (free X, Y), the map exists because the
%     positions make the linearisation observable. The paper treats co-estimating theta_base as out of
%     scope (footnote 1); we co-estimate. Precedent for linearising at zero: encoder paper Sec. 4.2
%     ("around the 0 point"), no reason given. TODOs: why Y = 0; disclose nominal parameters in W^b
%     versus detuned start; why model-based rather than data-based init (candidate: cost, Sec. 4.4);
%     Kaiming versus Xavier (todo only).
% 10. Encoder estimates the full state incl. measured positions (SUBNET). TODO: mention Kessels Remark 5.3?
% 11. Keep encoder history, scored window and validation horizon apart; values in Setup. The
%     joint-estimation sensitivity horizon belongs to Section IV.
% III-C Closed-loop identification
% 12. Known-controller, reference-driven output-error identification, inspired by Kessels (Eqs. 5.12,
%     5.13; Remark 5.6 open-loop attempt failed). Related to the indirect approach (Forssell and Ljung
%     Sec. 5.2) as a linear-theory motivation, not a guarantee; direct identification is not called
%     inherently biased. Free X and Y axes: the SUBNET consistency argument that Hoekstra inherits
%     (2026 Sec. 5.5) assumes an incrementally exponentially output stable data-generating system
%     (Beintema 2023 Condition 1); the open-loop gantry is not (free integrators), the loop with the
%     stabilising controller can be. Precondition, not a proof, for our nonlinear model. Same principle
%     as Kessels (closed-loop simulation of an unstable open-loop system), scoped to marginal stability.
% 13. Our formulation (ours, "we"): u_hat = u + C x^c + D(y - y_hat). Derivation by subtracting the
%     machine's controller equations, equal to the shared-r, shared-u_ff loop of the figure under a
%     compatible controller, signals and initial state (assumption; whether the data meet it is a Setup
%     or Results check: 4 kHz controller on recorded r - y versus recorded feedback force).
%     Main consequence: x^c_tau = 0 per window, no controller initialisation; consistent with the
%     encoder (model state and model controller both start on the machine's condition at tau). Contrast:
%     Kessels' direct form needs the machine controller state (Remark 5.4). Error before tau is not
%     carried in; the free-run validation does carry it.
% 14. Step order from no plant feedthrough: no algebraic loop, no added delay.
% 15. Every window sample scored (as Eq. 22); theta_base, theta_aug, eta joint, controller fixed.
%     Joint estimation framing (Dirk, 2026-10-06): the goal is to estimate theta_base jointly, in the ten
%     combinations. Obstacle from Jan's work: the split is not unique (2026 Sec. 5.2) and approximately
%     initialised parameters stay near their start (EJC Sec. 4.3, 2026 Sec. 6.3). Handover: Section IV
%     extends Jan's method with orthogonality so that the estimate of theta_base is UNIQUE (this
%     strength, not "recovers the true parameters"). State what interpretable means: the estimated
%     parameters keep their physical meaning; cite Gyorok 2026 Sec. 3 ("With non-unique theta_b, the
%     baseline model could lose its physical meaning"). Operational form only; the general definition
%     is the Introduction's (todo there). "Negation" and Kessels' definition (Sec. 5.3.1.3) belong to
%     Section IV; his Sec. 6.2 recommendation to the Introduction gap. Do NOT mention Jan's parameter regularisation anywhere.
%     Do NOT say explicitly that Section III's formulation is the unconstrained comparison arm; let
%     Section IV extend it naturally.
% Setup pattern for hyperparameters (2026 Sec. 7): honest basis per value (physical insight and
%     trial-and-error, line-search, inherited from earlier results).
% 16. Closed-loop fit can mask plant error; plant-domain checks and controller transfer in Results.
% Notation: u = P u_act as in Sections II and IV; no "(eom)" shorthand; controller matrices distinct from A_c(Y).
%
% PROPOSED ADDITIONS (writer, sandbox cycle 01, 2026-10-09; not yet approved by Dirk; derived from the
%   README ownership row for III and source-map rows III-C and III-D)
% 17. III-C also owns the training coordinates of theta_base (moved from Section II-C, eq:mdelta_map),
%     the optimiser (Adam, then an L-BFGS polish on a fixed batch of windows as in Drenth et al.) and
%     checkpoint selection (closed-loop free run over whole validation records, scored in metres). Values
%     stay in Experiment design. The polish acceptance meter compares against the true parameters (oracle,
%     sources.md trap 6): todo.
% 18. III-D PS2 (this work, motivated by Hefny 2015 two-stage regression and Downey 2017 2SR-then-BPTT):
%     problem = zero output layer leaves the added states without influence and without gradient; S1 fits
%     the encoder's added-state part and a temporary head to the current model's closed-loop residuals
%     over the next M = n+1 samples; S2 fits the deployed g_aug once to frozen encoder pairs; S3 continues
%     with the unchanged criterion. Deployed model, criterion and selection unchanged; switch rules use
%     training records only. Added states stay latent. Models: M_U (zero start) and M_PS2.
\section{Dynamic LPV-LFR Augmentation}\label{sec:augmentation}
\ifdraft
\begin{itemize}
\item The section extends the baseline with learned dynamics and identifies the result from closed-loop records; it names its four subsections (style profile)
\item Contribution share: the realisation on the self-scheduled gantry baseline with joint estimation of $\theta_{\mathrm{base}}$, its identification in closed loop with the known controller, and the initialisation of the added states (MUST ESTABLISH, item 18)
\end{itemize}
\fi

This section augments the baseline of Section~\ref{sec:lfr} with learned
dynamics and identifies the result from records measured under feedback.
Section~\ref{sec:aug_structure} defines the augmented model,
Section~\ref{sec:aug_encoder} the initial state of each training window,
Section~\ref{sec:aug_closed_loop} the closed-loop identification criterion and
its optimisation, and Section~\ref{sec:aug_ps2} the initialisation of the
additional states.
This work contributes the realisation on the self-scheduled gantry baseline,
with $\theta_{\mathrm{base}}$ estimated jointly in the combinations of
Section~\ref{sec:params}.
It further contributes the identification in closed loop with the known
controller and the initialisation of the additional states.
\todo{Missing: the Introduction's contribution list does not name the added-state initialisation (PS2). Candidate: extend the second contribution with ``and a data-driven initialisation of the additional states'' (D-236, D-240: PS2 is part of the reported method).}
\todo{VERIFY keys: \texttt{hoekstra2025lfr} = EJC 86 (2025) 101304,
\texttt{hoekstra2026lfr} = arXiv:2602.17297 (submitted to Automatica).
refs.bib and docs/references.md rows 46 and 130 now agree.}

\subsection{Augmented model}\label{sec:aug_structure}
\ifdraft
\begin{itemize}
\item If the omitted effects are dynamic, the six baseline states cannot represent them; refitting $\theta_{\mathrm{base}}$ or a static augmentation adds no states (EJC Sec.~2)
\item The model is Hoekstra's LFR augmentation structure with a fixed interconnection realising S-DP; its baseline block is the LFR of Section~II (EJC Sec.~3, 2026 Table~1)
\item Parallel because the orthogonal-by-construction method of Section~IV is defined for the additive structure (Gy\"or\"ok 2026); the baseline stays an explicit term
\item State augmentation only, output map not learned (?)
\item $f_{\mathrm{base}}$ is one RK4 step with the input held and $M(Y)$ per stage, so it stays self-scheduled; affine normalisation from the training records, unlike Hoekstra Eq.~(28)
\item The learned correction is added to the discrete-time transition after the step, as in the S-DP form (2026 Table~1)
\item One feedforward network with tanh hidden layers supplies $f_{\mathrm{aug}}$ and $g_{\mathrm{aug}}$; plain feedforward suffices for a parallel structure (EJC Sec.~4.3)
\item The correction acts on all physical rows, $X$ and $Y$ included; on these axes it meets no stiffness (?: why the position rows)
\item The network depends on $Y$ through the state but has no LPV structure, unlike Drenth Eq.~(5.2) (?)
\item The interconnection is acyclic, so it adds no algebraic loop; the baseline loop is well posed by Section~II
\item A zero output layer reproduces the baseline from the same physical initial state (2026 Eq.~(29)); additional states have no physical coordinates
\item The baseline/network split is not unique (2026 Sec.~5.2); interpretable means the estimated $\theta_{\mathrm{base}}$ keeps its physical meaning (Gy\"or\"ok Sec.~3); Section~IV makes the estimate unique
\end{itemize}
\fi

If the effects that the baseline omits are dynamic, its six states cannot
represent them.
Refitting $\theta_{\mathrm{base}}$ adds no states.
Neither does a static augmentation, which corrects the state update of the
baseline without extending its state~\cite[Sec.~2]{hoekstra2025lfr}.
The augmentation therefore needs states of its own.

We use the LFR-based augmentation structure of Hoekstra et
al.~\cite{hoekstra2025lfr,hoekstra2026lfr}, in which a fixed interconnection
connects a baseline block and a learning
block~\cite[Sec.~3.2]{hoekstra2025lfr}.
The baseline block is the self-scheduled LFR of Section~\ref{sec:lfr}, so the
augmented model is an LFR whose baseline block is itself an LFR in $Y$.
The interconnection realises the dynamic-parallel (S-DP)
structure~\cite[Table~1]{hoekstra2026lfr},
\begin{equation}
\begin{aligned}
 \tilde x_{k+1}
 &=f_{\mathrm{base}}(\theta_{\mathrm{base}},\tilde x_k,u_k)
   +f_{\mathrm{aug}}(\theta_{\mathrm{aug}},\hat x_k,u_k),\\
 \bar x_{k+1}
 &=g_{\mathrm{aug}}(\theta_{\mathrm{aug}},\hat x_k,u_k),\\
 \hat y_k&=h_{\mathrm{base}}(\tilde x_k),
\end{aligned}
\label{eq:aug_transition}
\end{equation}
where $\tilde x_k=\col(q_k,\dot q_k)\in\R^6$ is the physical state of
Section~\ref{sec:lfr}, $\bar x_k\in\R^{n_{\bar x}}$ the additional state,
$\hat x_k=\col(\tilde x_k,\bar x_k)$, $u_k=Pu_{\mathrm{act},k}\in\R^3$ the generalised
force of \eqref{eq:Ptransform}, and $h_{\mathrm{base}}(\tilde x)=P^\top q$ the
measured positions $\qact$.
We choose the parallel interconnection because the orthogonal-by-construction
parametrisation that Section~\ref{sec:obc} extends is defined for this
additive structure~\cite{gyorok2026obc}.
It also keeps the baseline as an explicit term of the state update.
The output map is not augmented.
\todo{Missing: reason for state augmentation only ($h_{\mathrm{aug}}=0$). Candidate: $h_{\mathrm{base}}$ is the exact kinematic map of the measured positions, the premise under which Kessels initialises positions from the output~\cite[Remark~5.3]{kessels2025ai}; an output augmentation depending on the input would also close an algebraic loop with the controller~\cite[Remark~5.1]{kessels2025ai}.}

The baseline transition $f_{\mathrm{base}}$ is one fourth-order Runge-Kutta
(RK4) step of the equations of motion \eqref{eq:eom} over the sample period
$T_s$, with the input held over the step.
Each stage evaluates $M(Y)$ at its own state, so $f_{\mathrm{base}}$ remains
self-scheduled.
As in the S-DP form, $f_{\mathrm{aug}}$ is added to the discrete-time
transition, after the step.
All signals are normalised with statistics of the training records,
\begin{equation}
 \tilde x=T_x\bigl(\col(q,\dot q)-\mu_x\bigr),\qquad
 T_x=\operatorname{diag}(\sigma_x)^{-1},
 \label{eq:aug_norm}
\end{equation}
where $\mu_x,\sigma_x\in\R^6$ are the mean and standard deviation of each
state channel.
The actuator forces $\uact$ and positions $\qact$ are normalised in the same
way.
Hoekstra et al.\ only scale the signals and take $\sigma_x$ from a simulation
of the baseline~\cite[Sec.~5.3]{hoekstra2026lfr}.
Here the state statistics come from the measured positions $q=P^{-\top}\qact$
and their first differences, since no state is measured.

One multilayer perceptron with tanh hidden layers and a linear output layer
supplies both learned functions,
\begin{equation}
\begin{aligned}
 \begin{bmatrix}
  f_{\mathrm{aug}}(\theta_{\mathrm{aug}},\hat x_k,u_k)\\
  g_{\mathrm{aug}}(\theta_{\mathrm{aug}},\hat x_k,u_k)
 \end{bmatrix}
 &=\phi_{\mathrm{aug}}(\theta_{\mathrm{aug}},\hat x_k,u_k)\\
 &=W_L\,\sigma_{L-1}(\hat x_k,u_k)+b_L,
\end{aligned}
 \label{eq:aug_network}
\end{equation}
where $\sigma_{L-1}$ is the output of the last hidden layer, $W_L$ and $b_L$
are the weight and bias of the output layer, $\theta_{\mathrm{aug}}$ collects
all weights and biases, and the first six rows form $f_{\mathrm{aug}}$
(Fig.~\ref{fig:aug_structure}).
A plain feedforward network suffices because the baseline is augmented
additively, whereas a series structure needs a residual network to reproduce
the baseline at initialisation~\cite[Sec.~4.3]{hoekstra2025lfr}.
As in the dynamic-parallel structure, $f_{\mathrm{aug}}$ corrects every
physical row~\cite[Sec.~2]{hoekstra2025lfr}, so the learned component can act
on the $X$ and $Y$ axes.
On these axes a correction meets no stiffness
(\eqref{eq:damping_stiffness_matrices}), so a persistent correction
accumulates in position.
\todo{Missing: why the position rows are corrected rather than kept as the integrals of the velocities. Context: Hoekstra et al.\ detach the free integrators of their vehicle and learn only the input-to-velocity map, which needs measured velocities~\cite[Sec.~7]{hoekstra2026lfr}; here only positions are measured. D-103 requires only that $X$ and $Y$ are reachable.}
The network depends on $Y$ through $\hat x$, so it can represent
position-dependent corrections, but unlike the augmented LPV-LFR of
Drenth~\cite[Eq.~(5.2)]{drenth2025thesis} it has no LPV structure.
\todo{Missing: reason for an unrestricted network rather than an LPV-structured augmentation. The research plan (Aspect~2) proposed the interconnection in LPV-LFR form, and D-017 (2026-03-19) still requires a $\Delta^a(Y)$ block; no later decision supersedes it, while the code uses an unrestricted MLP. Candidate: keep the plain S-DP network of Hoekstra et al.\ so that Section~IV's construction applies to one additive term (D-003, D-222).}
No learned function feeds the baseline block or itself within a step, so the
interconnection graph is acyclic and adds no algebraic
loop~\cite[Thm.~1]{hoekstra2025lfr}.
The scheduling loop of the baseline is well posed by
Section~\ref{sec:lfr}.

\begin{figure}[t]
 \centering
 \resizebox{\columnwidth}{!}{\input{figures/tex/augmentation_structure}}
 \caption{Augmented model. The physics step $f_{\mathrm{base}}$ and the
 network $\phi_{\mathrm{aug}}$ act in parallel on $\hat x_k$;
 $\phi_{\mathrm{aug}}$ supplies $f_{\mathrm{aug}}$ and $g_{\mathrm{aug}}$,
 and the output map is not learned ($h_{\mathrm{aug}}=0$). The encoder
 $\psi$ sets $\hat x$ at each window start. During training, $u_k$ is the
 closed-loop input $\hat u_k$.
 \todo{No setpoint generator currently for the inputs. to cramped together. We could make this a two column figure.}
 }
 \label{fig:aug_structure}
\end{figure}

Training starts from the baseline.
The output layer of the network is initialised at zero,
\begin{equation}
 W_L=0,\quad b_L=0
 \quad\Longrightarrow\quad
 f_{\mathrm{aug}}=0,\quad g_{\mathrm{aug}}=0,
 \label{eq:aug_zero_init}
\end{equation}
so that from the same physical initial state the augmented model reproduces
the baseline~\cite[Eq.~(29)]{hoekstra2026lfr}.
The additional states are zero after the first step and reach the output only
through $f_{\mathrm{aug}}$.
They have no assigned physical meaning: an invertible transformation of
$\bar x$, absorbed by the network, leaves the input-output behaviour
unchanged.
Nothing in the method ties them to a particular subsystem of the
data-generating system.
\todo{Missing: attribution of the all-zero output layer. It follows the authors' public dynamic-parallel code (D-237); \cite[Sec.~5.4.3]{hoekstra2026lfr} also allows a random linear part. Ask Jan which produced the reported S-DP results.}

From this start, $f_{\mathrm{aug}}$ can also learn changes that a different
$\theta_{\mathrm{base}}$ would produce, so the split between baseline and
network is not unique~\cite[Sec.~5.2]{hoekstra2026lfr}.
In the experiments of Hoekstra et al., approximately initialised physical
parameters stay close to their initial values, while the network learns
dynamics the baseline could represent~\cite[Sec.~6.3]{hoekstra2026lfr}.
With a non-unique $\theta_{\mathrm{base}}$, the baseline can lose its physical
meaning~\cite[Sec.~3]{gyorok2026obc}.
Here, the estimated combinations $\theta_{\mathrm{base}}$ are interpretable if
they keep that meaning, which requires a unique estimate.
Section~\ref{sec:obc} extends the method so that this estimate is unique.

\subsection{Encoder and initial state}\label{sec:aug_encoder}
\ifdraft
\begin{itemize}
\item Each window starts from a state of which only the positions are measured; an encoder estimates it from the input-output history (Beintema)
\item Windows starting from an estimated state keep the cost independent of the record length and stabilise the optimisation (Beintema Sec.~3.1)
\item The encoder estimates all states, the measured positions included (?: Kessels Remark~5.3)
\item The encoder is a linear map plus a nonlinear correction (encoder paper Eq.~(8))
\item $W^b$ starts from the reconstructability map of the baseline linearised at rest, $Y=0$ (encoder paper Eqs.~(15) to (17), (30) to (32)); computed, not fitted as in 2026 Eq.~(30)
\item The paper assumes a stable baseline; ours is marginally stable, but the positions make the linearisation observable
\item The initialisation improves starting values and convergence speed, not final accuracy (encoder paper Sec.~4.4)
\item $W^a$ starts from a Xavier draw, as in 2026 Eq.~(31) (resolved: thesis runs set \texttt{xavier\_uniform\_gain1}, D-239)
\item Open: why $Y=0$, nominal parameters in $W^b$, model-based versus data-based initialisation (?)
\end{itemize}
\fi

Each training window starts from a state of which only the positions are
measured.
Following Beintema et al.~\cite{beintema2023subnet}, an encoder estimates this
state from the input and output history,
\begin{equation}
\begin{aligned}
 \hat x_{\tau|\tau}
 &=\psi(\theta_{\mathrm{enc}},\xi_\tau),\\
 \xi_\tau&=\col\bigl(y_{\tau-n},\dots,y_\tau,\,u_{\tau-n},\dots,u_\tau\bigr),
\end{aligned}
 \label{eq:aug_encoder}
\end{equation}
where $\tau$ is the window start, $n$ the encoder history,
$\theta_{\mathrm{enc}}$ the encoder parameters, and $y=\qact$ and $u$ the
recorded output and input.
Simulating short windows from estimated states keeps the cost of a gradient
step independent of the record length and makes the optimisation more
stable~\cite[Sec.~3]{beintema2023subnet}.
The encoder estimates all states, the measured positions included.
\todo{Missing: whether to mention the alternative of initialising the positions from the measurement~\cite[Remark~5.3]{kessels2025ai}.}
\todo{Align the figure's past-only encoder-history labels with
\eqref{eq:aug_encoder}, which also uses the current sample.}

The encoder is a linear map with a nonlinear
correction~\cite[Eq.~(8)]{hoekstra2026encoder},
\begin{equation}
 \psi(\theta_{\mathrm{enc}},\xi_\tau)
 =\begin{bmatrix}W^b\\ W^a\end{bmatrix}\xi_\tau+\tilde\psi(\xi_\tau),
 \label{eq:aug_enc_lin}
\end{equation}
where the rows $W^b$ and $W^a$ give the physical and the additional states,
and $\tilde\psi$ is a multilayer perceptron whose output layer starts at zero.
$W^a$ has no baseline counterpart and starts from a Xavier draw, as for the
additional-state encoder of~\cite[Eq.~(31)]{hoekstra2026lfr}.
\todo{VERIFY bib: \texttt{hoekstra2026encoder} = arXiv:2602.13108v1
(13 Feb 2026), Hoekstra, Gy\"or\"ok, T\'oth, Schoukens; matches PDF p.~1.
Check for a published (IFAC) version before submission.}

The physical rows start from the baseline, following the model-based
initialisation of Hoekstra et al.~\cite[Sec.~3.1]{hoekstra2026encoder}.
For zero input, every rest state with $\Theta=0$ is an equilibrium of
\eqref{eq:eom}, because
$K$ in \eqref{eq:damping_stiffness_matrices} vanishes on $X$ and $Y$.
We linearise the baseline at rest at $Y=0$ and discretise it by zero-order
hold over $T_s$, which gives $(A_d,B_d,C_d)$ in normalised coordinates.
The reconstructability map of this model~\cite[Eqs.~(15) to (17)]{hoekstra2026encoder}
gives
\begin{equation}
 W^b=\begin{bmatrix}
  A_d^nO_n^{+} & \;r_n-A_d^nO_n^{+}T_n
 \end{bmatrix},
 \label{eq:aug_wb}
\end{equation}
where $O_n$ is the observability matrix of $(A_d,C_d)$ over $n+1$ samples,
$T_n$ the Toeplitz matrix of its Markov parameters, $r_n$ the input-to-state
map over the same samples, and the two blocks act on the output and input
parts of $\xi_\tau$.
The encoder paper assumes a stable baseline.
Ours is marginally stable, but the measured positions make $(A_d,C_d)$
observable, so $O_n$ has full column rank and the map exists.
Unlike~\cite[Eq.~(30)]{hoekstra2026lfr}, which fits the baseline encoder to
simulated baseline states, we compute the map directly.
The model-based initialisation gives better starting values and faster
convergence than a random encoder, at comparable final
accuracy~\cite[Sec.~4.4]{hoekstra2026encoder}.
The encoder is trained jointly with the model, so the linearisation only sets
its starting value.
\todo{Missing: why linearise at $Y=0$. The encoder paper also linearises ``around the 0 point''~\cite[Sec.~4.2]{hoekstra2026encoder} but gives no reason.}
\todo{Missing: disclosure. $W^b$ is built from the nominal parameters (\texttt{gantry\_linearize\_and\_discretize} in \texttt{GD/model.py::build\_model}), which in simulation equal the true ones, also when $\theta_{\mathrm{base}}$ starts detuned; the encoder paper leaves co-estimation out of scope~\cite[footnote~1]{hoekstra2026encoder}. Candidate: state it as prior information of the grey-box model and weigh it in the black-box comparison of Section~\ref{sec:setup} (D-243 gives the black box no physical knowledge).}
\todo{Missing: why the model-based rather than the data-based initialisation~\cite[Sec.~3.2]{hoekstra2026encoder}. Candidate: negligible computation~\cite[Sec.~4.4]{hoekstra2026encoder}.}

\subsection{Closed-loop identification}\label{sec:aug_closed_loop}
\ifdraft
\begin{itemize}
\item The records are measured under feedback and the model must predict the machine under feedback, so we fit the closed-loop response, not a rollout driven by the recorded input (2026 Eq.~(22))
\item Inspired by Kessels (Eqs.~(5.12), (5.13); Remark~5.6); for linear systems the indirect approach, where an output-error criterion remains consistent (Forssell and Ljung); a motivation, not a guarantee
\item $X$ and $Y$ are free, so the open-loop gantry violates the stability condition behind SUBNET consistency (Beintema Cond.~1); the loop with the controller can satisfy it
\item Both loops share $r$, $u^{\mathrm{ff}}$ and the controller; subtracting them gives our residual form, exact under a compatible controller, signals and initial state (check in Setup)
\item $x^{c}_\tau=0$ is the true controller state that Kessels' Remark~5.4 supplies, obtained without reading or reconstructing it; it matches the encoder
\item No plant feedthrough fixes the step order: no algebraic loop, no added delay
\item $\theta_{\mathrm{base}}$, $\theta_{\mathrm{aug}}$ and $\theta_{\mathrm{enc}}$ minimise the closed-loop output error over every window sample, controller fixed; the window spans the closed-loop memory (D-220)
\item Moved from Section~II-C: state and justify the training coordinates of $\theta_{\mathrm{base}}$ here. Nine positive combinations in logarithmic coordinates; the signed $m_\Delta$ through \eqref{eq:mdelta_map}. Justification: every $\xi_\Delta$ satisfies \eqref{eq:lfr_admissibility}, so $M(Y)\succ0$ and the LFR stays well posed at every optimiser iterate without clamp or penalty (D-229); constraint by parameterization as in \cite{sutanto2020encoding}; it guarantees positive-definite inertia, not positive raw masses (positive $m_1,m_2$ need $|m_\Delta|<m_\Sigma-m_b$, with $m_b$ unidentifiable); $\eta_0$, $s$ in Appendix~\ref{app:lfr-wellposed}, $\rho$ in Setup. Notation clash resolved: encoder parameters are $\theta_{\mathrm{enc}}$ and controller states $x^{\mathrm{fb}}$, so $\eta_0$, $s$, $\xi_\Delta$ are free
\item Adam on minibatches of windows, then an L-BFGS polish on one fixed set of windows, kept only if the validation score improves (Drenth; D-171, D-174)
\item The checkpoint with the lowest closed-loop free-run RMS on the validation records is selected; the free run carries model error over the whole record (D-171, D-172, D-198)
\item The horizons are separate quantities with values in Setup; a closed-loop fit can hide plant error (Results)
\item Model names: $\mathcal M_{\mathrm{U}}$ is this model and criterion from the zero start (?: notation)
\end{itemize}
\fi

The records are measured under feedback, and the model must predict the
machine under feedback: its response to a reference, a feedforward and a
known controller.
We therefore identify the model from this closed-loop response, rather than
from a simulation driven by the recorded input as
in~\cite[Eq.~(22)]{hoekstra2026lfr}.
Kessels trains augmented models in this way, simulating each window in closed
loop with the known controller~\cite[Eqs.~(5.12), (5.13)]{kessels2025ai}.
He also reports that identifying the open-loop model from his closed-loop data
failed~\cite[Remark~5.6]{kessels2025ai}.
For linear systems, fitting the closed loop from the reference with a known
controller is the indirect approach to closed-loop
identification~\cite[Sec.~5.2]{forssell1999revisited}.
The reference is uncorrelated with the noise, so an output-error criterion
remains consistent under the conditions given there.
For our nonlinear model, this is a motivation, not a guarantee.

The free axes of the baseline give a second reason.
$X$ and $Y$ have no stiffness in \eqref{eq:damping_stiffness_matrices}, so the
open-loop gantry is marginally stable:
a difference in initial velocity leaves a permanent position offset.
The consistency of the identification method~\cite[Sec.~5.5]{hoekstra2026lfr}
rests on that of SUBNET, which assumes an incrementally exponentially output
stable data-generating system~\cite[Cond.~1]{beintema2023subnet}.
The open-loop gantry does not satisfy this condition, whereas the loop with
the stabilising controller can.
Kessels uses the same principle to update models of systems that are unstable
in open loop~\cite[Ch.~5]{kessels2025ai}.
For our nonlinear model, this makes the condition attainable but does not
prove consistency.
\todo{Missing: evidence on the thesis data that open-loop rollouts accumulate position error on $X$ and $Y$; measure and report in Results (THESIS-RESULTS Sec.~5, ``Section III-C evidence'').}

Both loops are driven by the same reference $r$ and feedforward
$u^{\mathrm{ff}}$ through the same controller, and differ only in the output
they feed back (Fig.~\ref{fig:aug_closed_loop}).
The controller of each record is known, linear and time invariant, with
state-space matrices $(A_{\mathrm{fb}},B_{\mathrm{fb}},C_{\mathrm{fb}},D_{\mathrm{fb}})$
mapping the position error in actuator coordinates to actuator forces.
With controller state $x^{\mathrm{fb}}$ on the machine and
$\hat x^{\mathrm{fb}}$ in the model, the loops are
\begin{equation}
\begin{aligned}
 x^{\mathrm{fb}}_{k+1}&=A_{\mathrm{fb}}x^{\mathrm{fb}}_k+B_{\mathrm{fb}}(r_k-y_k),\\
 u_{\mathrm{act},k}&=C_{\mathrm{fb}}x^{\mathrm{fb}}_k+D_{\mathrm{fb}}(r_k-y_k)+u^{\mathrm{ff}}_k,\\
 \hat x^{\mathrm{fb}}_{k+1}&=A_{\mathrm{fb}}\hat x^{\mathrm{fb}}_k+B_{\mathrm{fb}}(r_k-\hat y_k),\\
 \hat u_{\mathrm{act},k}&=C_{\mathrm{fb}}\hat x^{\mathrm{fb}}_k+D_{\mathrm{fb}}(r_k-\hat y_k)+u^{\mathrm{ff}}_k.
\end{aligned}
\label{eq:aug_direct_loop}
\end{equation}
We implement the model loop by subtracting the machine's equations from the
model's.
This removes $r$ and $u^{\mathrm{ff}}$ and gives
\begin{equation}
\begin{aligned}
 x^{c}_{k+1}&=A_{\mathrm{fb}}x^{c}_k+B_{\mathrm{fb}}(y_k-\hat y_k),\qquad x^{c}_\tau=0,\\
 \hat u_{\mathrm{act},k}&=u_{\mathrm{act},k}+C_{\mathrm{fb}}x^{c}_k+D_{\mathrm{fb}}(y_k-\hat y_k),
\end{aligned}
\label{eq:aug_residual_loop}
\end{equation}
where $x^{c}=\hat x^{\mathrm{fb}}-x^{\mathrm{fb}}$.
The model thus receives the recorded input plus the controller's response to
its own output error, and $\hat u_k=P\hat u_{\mathrm{act},k}$ replaces $u_k$
in \eqref{eq:aug_transition}.
The two forms are equal when the recorded input was produced by this
controller from the same output samples and initial state.
\todo{Check for Setup: the training loop uses the controller re-discretised at the model rate, while the records were generated at a higher rate. Run the re-discretised controller on the recorded $r-y$ and compare its feedback force with the recorded one. A rollout with $\hat y=y$ gives $\hat u=u$ by construction and checks nothing.}

\begin{figure}[t]
 \centering
 \includegraphics[width=\columnwidth]{closed_loop_same_reference.pdf}
 \caption{Recorded plant loop (top) and augmented-model loop (bottom),
 driven by one shared reference $r_k$ and feedforward $u_k^{\mathrm{ff}}$,
 with the same feedback controller. The model output is evaluated
 before the input advances its state to the next sample.}
 \label{fig:aug_closed_loop}
\end{figure}

Because $x^{c}$ is a difference of controller states, $x^{c}_\tau=0$ places
the model's controller in the machine's controller state,
$\hat x^{\mathrm{fb}}_\tau=x^{\mathrm{fb}}_\tau$, just as the encoder places
the model state on the recorded trajectory.
This is the initialisation that Kessels obtains by reading the machine's
controller state or reconstructing it from the reference, the output and the
controller~\cite[Remark~5.4]{kessels2025ai}.
The residual form needs neither.
Model error made before $\tau$ is not carried into the window.
The augmented model has no feedthrough, so each sample is evaluated in a fixed
order: $\hat y_k$ from $\tilde x_k$, then $\hat u_{\mathrm{act},k}$ from
\eqref{eq:aug_residual_loop}, then the model and controller states advance.
The controller feedthrough $D_{\mathrm{fb}}$ therefore closes no algebraic
loop, and the loop adds no delay.
\todo{Resolve in docs: D-142 states that $x^c=0$ \emph{is} Kessels' Remark~5.4, while \texttt{FS/closed\_loop.py::closed\_loop\_rollout} says it is not. The text above takes the reading that both hold: $x^c_\tau=0$ yields the true controller state Remark~5.4 supplies, without its reconstruction. Candidate: amend the code comment.}

The physical parameters are trained in coordinates that keep the baseline
well posed.
The nine positive combinations of \eqref{eq:phi} are trained as logarithms.
The signed $m_\Delta$ is trained through
\begin{equation}
  m_\Delta=\rho\,\frac{2}{L_b}\sqrt{m_\Sigma J_{\mathrm{eff}}}\,\tanh(\eta_0+s\,\xi_\Delta),
  \label{eq:mdelta_map}
\end{equation}
where $\xi_\Delta\in\R$ is the trained coordinate, $\rho<1$ a margin, and
$\eta_0$ and $s$ fix the start and the local scale
(Appendix~\ref{app:lfr-wellposed}).
Every $\xi_\Delta$ satisfies \eqref{eq:lfr_admissibility}, so $M(Y)\succ0$
and the LFR stays well posed at every iterate without a clamp or a penalty.
As in~\cite{sutanto2020encoding}, the constraint is enforced by the
parametrisation.
The map guarantees positive-definite inertia, not positive raw masses.
\todo{VERIFY citation: \cite{sutanto2020encoding} has no row in docs/references.md and its passage was not read in this cycle (README source map II-C).}

Over the window starts $\mathcal T$ and the window length $n_f$, we identify
the model as
\begin{equation}
\begin{aligned}
 \hat\vartheta&=\arg\min_{\vartheta}V(\vartheta),\\
 V(\vartheta)&=\frac{1}{|\mathcal T|\,n_f\,n_y}
 \sum_{\tau\in\mathcal T}\sum_{\ell=0}^{n_f-1}
 \norm{y_{\tau+\ell}-\hat y_{\tau+\ell|\tau}}_2^2,\\
 &\text{s.t.}\quad \hat x_{\tau|\tau}=\psi(\theta_{\mathrm{enc}},\xi_\tau),\quad x^c_\tau=0,\\
 &\phantom{\text{s.t.}\quad}\text{\eqref{eq:aug_transition} with $u=\hat u$, and \eqref{eq:aug_residual_loop}},
\end{aligned}
\label{eq:aug_loss}
\end{equation}
where $\vartheta=(\theta_{\mathrm{base}},\theta_{\mathrm{aug}},\theta_{\mathrm{enc}})$,
$\theta_{\mathrm{base}}$ enters through the coordinates above, $n_y=3$, the
outputs are normalised, and $\hat y_{\tau+\ell|\tau}$ is the output $\ell$
samples after $\tau$ simulated from $\hat x_{\tau|\tau}$.
As in~\cite[Eq.~(22)]{hoekstra2026lfr}, every sample of a window is scored.
Both initial conditions are set by the window start, so no sample needs to be
discarded as a transient.
$\theta_{\mathrm{base}}$ is estimated jointly with the network and the
encoder, while the controller stays fixed.
Kessels' criterion, in contrast, leaves the first-principles parameters
fixed~\cite[Eq.~(5.12)]{kessels2025ai}.
Since $x^c$ restarts at every window, $n_fT_s$ must span the time a residual
force stays visible in the closed-loop output.
The poles of the loop with the controller bound that time.
\todo{Missing: the window bound is computed from the linearised truth with the controller (D-220, option A), i.e.\ from model matrices of the data-generating system. Candidate: state in Section~\ref{sec:setup} that the value is derived from the loop poles and confirmed by the measured closed-loop FRF (option B), not tuned.}

As in Drenth et al.~\cite{drenth2025lpvlfr}, the criterion is minimised with
Adam and then refined with L-BFGS.
Adam uses minibatches of windows.
L-BFGS runs a fixed number of iterations with a strong-Wolfe line search on
one fixed set of windows.
Unlike there, L-BFGS runs unconstrained, since \eqref{eq:mdelta_map} already
keeps every iterate admissible.
The model that is reported is selected on the validation records by a
closed-loop free run: from $\hat x_{k_0|k_0}$ and $x^c_{k_0}=0$, each record
$i$ is simulated through \eqref{eq:aug_transition} and
\eqref{eq:aug_residual_loop} to its end, and scored by
\begin{equation}
\begin{aligned}
 J_{\mathrm{val}}
 &=\frac{\sum_i N_i\,e_i}{\sum_i N_i},\\
 e_i&=\Bigl(\frac{1}{(N_i-k_0)\,n_y}\sum_{k=k_0}^{N_i-1}
 \norm{y_k-\hat y_{k|k_0}}_2^2\Bigr)^{1/2},
\end{aligned}
 \label{eq:aug_val}
\end{equation}
where $N_i$ is the record length, $k_0=n$ the first sample with a complete
encoder history, and $y$ is in metres.
Unlike a training window, the free run carries model error over the whole
record.
During Adam, the checkpoint with the lowest $J_{\mathrm{val}}$ is kept.
The polished parameters replace it only if they lower $J_{\mathrm{val}}$ and
no monitored parameter meter regresses.
The encoder history $n$, the window length $n_f$ and the free-run horizon are
separate quantities, whose values are given in Section~\ref{sec:setup}.
\todo{Missing: the polish acceptance meter \texttt{combo\_err} compares the combinations with the true parameters (\texttt{GD/training.py::\_gantry\_meters}), an oracle inside training. Candidate: drop that meter from acceptance, or state it as a simulation-only safeguard.}

We call the model \eqref{eq:aug_transition} with encoder
\eqref{eq:aug_encoder}, started from \eqref{eq:aug_zero_init} and identified
by \eqref{eq:aug_loss}, $\mathcal M_{\mathrm{U}}$.
A good closed-loop fit does not by itself show that the model is correct in
open loop, because the controller acts on the output error and can mask plant
error.
Section~\ref{sec:results} therefore also examines the model outside the
training loop.
\todo{Missing: model-name notation. The README writes $M_{\mathrm{U}}$, which clashes with the inertia matrix $M(Y)$. Candidate: $\mathcal M_{\mathrm{U}}$, $\mathcal M_{\mathrm{PS2}}$, $\mathcal M_{\mathrm{OBC}}$ (notation is Dirk's decision).}

\subsection{Initialisation of the additional states}\label{sec:aug_ps2}
\ifdraft
\begin{itemize}
\item From \eqref{eq:aug_zero_init}, every path from $\bar x$ to $\hat y$ passes through $W_L=0$: $W^a$ and the $g_{\mathrm{aug}}$ rows of $W_L$ start with zero gradient of $V$ (this work, PS2-METHOD Sec.~2.2)
\item $W^b$ gives the physical states a meaning from the baseline; no baseline exists for $\bar x$, so we give them one from data, before they influence the output (this work, D-236)
\item S1: a temporary head and the encoder's additional-state part regress the current model's closed-loop residuals over the next $M=n+1$ samples (this work; motivated by Hefny S1)
\item S2: once, the deployed $g_{\mathrm{aug}}$ is fitted to frozen encoder pairs, one step (this work; motivated by Hefny S2 and Downey Sec.~4.2)
\item S3: the head is discarded and training continues with $V$; model, criterion and selection unchanged (code, D-236)
\item Phase switches use training records only; a failed screen skips S2; thresholds in Setup (code, PS2-METHOD Sec.~4.5)
\item Limit: the added states stay latent; one-step fit does not guarantee multi-step propagation (PS2-METHOD Sec.~6.2.1)
\item Model name: $\mathcal M_{\mathrm{PS2}}$ (?: notation)
\end{itemize}
\fi

The zero output layer of \eqref{eq:aug_zero_init} leaves the additional states
without influence and without a learning signal.
Every path from $\bar x$ to $\hat y$ passes through $W_L$, so at the start
the gradient of $V$ is zero with respect to $W^a$ and to the rows of $W_L$
that produce $g_{\mathrm{aug}}$.
These parameters receive a gradient only once the rows producing
$f_{\mathrm{aug}}$ depend on $\bar x$, that is, through products of weights
that start at zero.
The encoder gives the physical states a meaning from the baseline through
$W^b$.
No baseline exists for $\bar x$, so we propose to give the additional states a
meaning from data before they reach the output.
The procedure, PS2, has two supervised stages during the first updates of the
training and changes only the starting values of $\theta_{\mathrm{enc}}$ and
$\theta_{\mathrm{aug}}$.

The first stage (S1) makes the encoded additional states predictive of what
the current model gets wrong.
For a window start $\tau$, the target is the closed-loop output residual of
the current model over the next $M=n+1$ samples, normalised by its root mean
square $\sigma_e$ over the minibatch $\mathcal B$,
\begin{equation}
\begin{aligned}
 e_\tau&=\sigma_e^{-1}\col\bigl(y_{\tau+\ell}-\hat y_{\tau+\ell|\tau}\bigr)_{\ell=0}^{M-1},\\
 \min_{\phi,\,\theta_{\mathrm{enc}}^{a}}&\;
 \frac{1}{|\mathcal B|\,M\,n_y}\sum_{\tau\in\mathcal B}
 \norm{h_\phi\bigl(\psi_a(\theta_{\mathrm{enc}},\xi_\tau)\bigr)-e_\tau}_2^2,
\end{aligned}
\label{eq:ps2_s1}
\end{equation}
where $\hat y_{\tau+\ell|\tau}$ is predicted as in \eqref{eq:aug_loss} and
held fixed, $\psi_a\in\R^{n_{\bar x}}$ are the additional-state rows of the
encoder, $\theta_{\mathrm{enc}}^{a}$ are $W^a$ and the rows of the output
layer of $\tilde\psi$ that produce $\psi_a$, and
$h_\phi:\R^{n_{\bar x}}\to\R^{Mn_y}$ is a temporary head with one tanh hidden
layer and an output layer that starts at zero.
One S1 step with its own Adam optimiser follows every update of $V$, on
windows of the training records except a held-out calibration subset.
S1 ends when the fraction of $\norm{e}^2$ that $h_\phi$ leaves unexplained on
the calibration windows stops decreasing, or after a fixed number of updates.

The second stage (S2) makes the deployed transition propagate these states.
Once, at the end of S1, encoder pairs one sample apart are drawn and frozen,
and the network is fitted by
\begin{equation}
 \min_{\theta_{\mathrm{aug}}}\;
 \sum_{k}\norm{g_{\mathrm{aug}}\bigl(\theta_{\mathrm{aug}},\hat x_{k|k},u_k\bigr)
 -\psi_a(\theta_{\mathrm{enc}},\xi_{k+1})}_2^2,
 \label{eq:ps2_s2}
\end{equation}
with $\hat x_{k|k}=\psi(\theta_{\mathrm{enc}},\xi_k)$, normalised by the
variance of the targets.
The rows of $W_L$ that produce $f_{\mathrm{aug}}$ receive no gradient from
\eqref{eq:ps2_s2} and stay unchanged.
The hidden layers, shared with $f_{\mathrm{aug}}$, do move.
S2 runs until its residual on calibration pairs stops decreasing and keeps
the weights with the lowest residual.
If at the end of S1 the head still leaves most of $\norm{e}^2$ unexplained,
S2 is skipped.
The switching rules use training records only.
Their thresholds are listed in Section~\ref{sec:setup}.

After S2 the head is discarded (S3), and training continues with $V$,
the optimiser and the selection of Section~\ref{sec:aug_closed_loop}.
The deployed model, the criterion and the selection are those of
$\mathcal M_{\mathrm{U}}$.
We call the model obtained with PS2 $\mathcal M_{\mathrm{PS2}}$.
S1 follows the first stage of the two-stage regression of Hefny et al., which
regresses a predictive state on the history~\cite[Sec.~3]{hefny2015supervised}.
Unlike there, the target is the residual of the current closed-loop model and
the state is the encoder output.
S2 replaces their linear operator between stage-one outputs by a one-step fit
of the nonlinear $g_{\mathrm{aug}}$.
Initialising by two-stage regression and then refining by backpropagation
through time follows Downey et al.~\cite[Sec.~4.2]{downey2017psrnn}.
Neither stage ties $\bar x$ to a physical subsystem.
A one-step fit also does not guarantee that $g_{\mathrm{aug}}$ propagates the
states over many steps.
S3 refines the transition through $V$.
\todo{Missing: the S2 weight selection in \texttt{FS/ps2\_opening.py::PS2Opening.\_fit\_map} keeps the final weights when no calibration check improves on the start (PS2-METHOD Sec.~4.3, known defect). Candidate: state ``keeps the best weights, the start included'' after the code is fixed, or report the S2 residual of each run in Results.}
